% Options for packages loaded elsewhere
% Options for packages loaded elsewhere
\PassOptionsToPackage{unicode}{hyperref}
\PassOptionsToPackage{hyphens}{url}
\PassOptionsToPackage{dvipsnames,svgnames,x11names}{xcolor}
%
\documentclass[
  12pt,
  letterpaper,
  DIV=11,
  numbers=noendperiod]{scrartcl}
\usepackage{xcolor}
\usepackage{amsmath,amssymb}
\setcounter{secnumdepth}{5}
\usepackage{iftex}
\ifPDFTeX
  \usepackage[T1]{fontenc}
  \usepackage[utf8]{inputenc}
  \usepackage{textcomp} % provide euro and other symbols
\else % if luatex or xetex
  \usepackage{unicode-math} % this also loads fontspec
  \defaultfontfeatures{Scale=MatchLowercase}
  \defaultfontfeatures[\rmfamily]{Ligatures=TeX,Scale=1}
\fi
\usepackage{lmodern}
\ifPDFTeX\else
  % xetex/luatex font selection
  \setmainfont[]{Times New Roman}
\fi
% Use upquote if available, for straight quotes in verbatim environments
\IfFileExists{upquote.sty}{\usepackage{upquote}}{}
\IfFileExists{microtype.sty}{% use microtype if available
  \usepackage[]{microtype}
  \UseMicrotypeSet[protrusion]{basicmath} % disable protrusion for tt fonts
}{}
\makeatletter
\@ifundefined{KOMAClassName}{% if non-KOMA class
  \IfFileExists{parskip.sty}{%
    \usepackage{parskip}
  }{% else
    \setlength{\parindent}{0pt}
    \setlength{\parskip}{6pt plus 2pt minus 1pt}}
}{% if KOMA class
  \KOMAoptions{parskip=half}}
\makeatother
% Make \paragraph and \subparagraph free-standing
\makeatletter
\ifx\paragraph\undefined\else
  \let\oldparagraph\paragraph
  \renewcommand{\paragraph}{
    \@ifstar
      \xxxParagraphStar
      \xxxParagraphNoStar
  }
  \newcommand{\xxxParagraphStar}[1]{\oldparagraph*{#1}\mbox{}}
  \newcommand{\xxxParagraphNoStar}[1]{\oldparagraph{#1}\mbox{}}
\fi
\ifx\subparagraph\undefined\else
  \let\oldsubparagraph\subparagraph
  \renewcommand{\subparagraph}{
    \@ifstar
      \xxxSubParagraphStar
      \xxxSubParagraphNoStar
  }
  \newcommand{\xxxSubParagraphStar}[1]{\oldsubparagraph*{#1}\mbox{}}
  \newcommand{\xxxSubParagraphNoStar}[1]{\oldsubparagraph{#1}\mbox{}}
\fi
\makeatother

\usepackage{color}
\usepackage{fancyvrb}
\newcommand{\VerbBar}{|}
\newcommand{\VERB}{\Verb[commandchars=\\\{\}]}
\DefineVerbatimEnvironment{Highlighting}{Verbatim}{commandchars=\\\{\}}
% Add ',fontsize=\small' for more characters per line
\usepackage{framed}
\definecolor{shadecolor}{RGB}{241,243,245}
\newenvironment{Shaded}{\begin{snugshade}}{\end{snugshade}}
\newcommand{\AlertTok}[1]{\textcolor[rgb]{0.68,0.00,0.00}{#1}}
\newcommand{\AnnotationTok}[1]{\textcolor[rgb]{0.37,0.37,0.37}{#1}}
\newcommand{\AttributeTok}[1]{\textcolor[rgb]{0.40,0.45,0.13}{#1}}
\newcommand{\BaseNTok}[1]{\textcolor[rgb]{0.68,0.00,0.00}{#1}}
\newcommand{\BuiltInTok}[1]{\textcolor[rgb]{0.00,0.23,0.31}{#1}}
\newcommand{\CharTok}[1]{\textcolor[rgb]{0.13,0.47,0.30}{#1}}
\newcommand{\CommentTok}[1]{\textcolor[rgb]{0.37,0.37,0.37}{#1}}
\newcommand{\CommentVarTok}[1]{\textcolor[rgb]{0.37,0.37,0.37}{\textit{#1}}}
\newcommand{\ConstantTok}[1]{\textcolor[rgb]{0.56,0.35,0.01}{#1}}
\newcommand{\ControlFlowTok}[1]{\textcolor[rgb]{0.00,0.23,0.31}{\textbf{#1}}}
\newcommand{\DataTypeTok}[1]{\textcolor[rgb]{0.68,0.00,0.00}{#1}}
\newcommand{\DecValTok}[1]{\textcolor[rgb]{0.68,0.00,0.00}{#1}}
\newcommand{\DocumentationTok}[1]{\textcolor[rgb]{0.37,0.37,0.37}{\textit{#1}}}
\newcommand{\ErrorTok}[1]{\textcolor[rgb]{0.68,0.00,0.00}{#1}}
\newcommand{\ExtensionTok}[1]{\textcolor[rgb]{0.00,0.23,0.31}{#1}}
\newcommand{\FloatTok}[1]{\textcolor[rgb]{0.68,0.00,0.00}{#1}}
\newcommand{\FunctionTok}[1]{\textcolor[rgb]{0.28,0.35,0.67}{#1}}
\newcommand{\ImportTok}[1]{\textcolor[rgb]{0.00,0.46,0.62}{#1}}
\newcommand{\InformationTok}[1]{\textcolor[rgb]{0.37,0.37,0.37}{#1}}
\newcommand{\KeywordTok}[1]{\textcolor[rgb]{0.00,0.23,0.31}{\textbf{#1}}}
\newcommand{\NormalTok}[1]{\textcolor[rgb]{0.00,0.23,0.31}{#1}}
\newcommand{\OperatorTok}[1]{\textcolor[rgb]{0.37,0.37,0.37}{#1}}
\newcommand{\OtherTok}[1]{\textcolor[rgb]{0.00,0.23,0.31}{#1}}
\newcommand{\PreprocessorTok}[1]{\textcolor[rgb]{0.68,0.00,0.00}{#1}}
\newcommand{\RegionMarkerTok}[1]{\textcolor[rgb]{0.00,0.23,0.31}{#1}}
\newcommand{\SpecialCharTok}[1]{\textcolor[rgb]{0.37,0.37,0.37}{#1}}
\newcommand{\SpecialStringTok}[1]{\textcolor[rgb]{0.13,0.47,0.30}{#1}}
\newcommand{\StringTok}[1]{\textcolor[rgb]{0.13,0.47,0.30}{#1}}
\newcommand{\VariableTok}[1]{\textcolor[rgb]{0.07,0.07,0.07}{#1}}
\newcommand{\VerbatimStringTok}[1]{\textcolor[rgb]{0.13,0.47,0.30}{#1}}
\newcommand{\WarningTok}[1]{\textcolor[rgb]{0.37,0.37,0.37}{\textit{#1}}}

\usepackage{longtable,booktabs,array}
\newcounter{none} % for unnumbered tables
\usepackage{calc} % for calculating minipage widths
% Correct order of tables after \paragraph or \subparagraph
\usepackage{etoolbox}
\makeatletter
\patchcmd\longtable{\par}{\if@noskipsec\mbox{}\fi\par}{}{}
\makeatother
% Allow footnotes in longtable head/foot
\IfFileExists{footnotehyper.sty}{\usepackage{footnotehyper}}{\usepackage{footnote}}
\makesavenoteenv{longtable}
\usepackage{graphicx}
\makeatletter
\newsavebox\pandoc@box
\newcommand*\pandocbounded[1]{% scales image to fit in text height/width
  \sbox\pandoc@box{#1}%
  \Gscale@div\@tempa{\textheight}{\dimexpr\ht\pandoc@box+\dp\pandoc@box\relax}%
  \Gscale@div\@tempb{\linewidth}{\wd\pandoc@box}%
  \ifdim\@tempb\p@<\@tempa\p@\let\@tempa\@tempb\fi% select the smaller of both
  \ifdim\@tempa\p@<\p@\scalebox{\@tempa}{\usebox\pandoc@box}%
  \else\usebox{\pandoc@box}%
  \fi%
}
% Set default figure placement to htbp
\def\fps@figure{htbp}
\makeatother





\setlength{\emergencystretch}{3em} % prevent overfull lines

\providecommand{\tightlist}{%
  \setlength{\itemsep}{0pt}\setlength{\parskip}{0pt}}



 


\usepackage{tcolorbox}
\usepackage{amssymb}
\usepackage{yfonts}
\usepackage{bm}


\newtcolorbox{greybox}{
  colback=white,
  colframe=blue,
  coltext=black,
  boxsep=5pt,
  arc=4pt}
  
\newcommand{\sectionbreak}{\clearpage}

 
\newcommand{\ds}[4]{\sum_{{#1}=1}^{#3}\sum_{{#2}=1}^{#4}}
\newcommand{\us}[3]{\mathop{\sum\sum}_{1\leq{#2}<{#1}\leq{#3}}}

\newcommand{\ol}[1]{\overline{#1}}
\newcommand{\ul}[1]{\underline{#1}}

\newcommand{\amin}[1]{\mathop{\text{argmin}}_{#1}}
\newcommand{\amax}[1]{\mathop{\text{argmax}}_{#1}}

\newcommand{\ci}{\perp\!\!\!\perp}

\newcommand{\mc}[1]{\mathcal{#1}}
\newcommand{\mb}[1]{\mathbb{#1}}
\newcommand{\mf}[1]{\mathfrak{#1}}

\newcommand{\eps}{\epsilon}
\newcommand{\lbd}{\lambda}
\newcommand{\alp}{\alpha}
\newcommand{\df}{=:}
\newcommand{\am}[1]{\mathop{\text{argmin}}_{#1}}
\newcommand{\ls}[2]{\mathop{\sum\sum}_{#1}^{#2}}
\newcommand{\ijs}{\mathop{\sum\sum}_{1\leq i<j\leq n}}
\newcommand{\jis}{\mathop{\sum\sum}_{1\leq j<i\leq n}}
\newcommand{\sij}{\sum_{i=1}^n\sum_{j=1}^n}
	
\KOMAoption{captions}{tableheading}
\makeatletter
\@ifpackageloaded{caption}{}{\usepackage{caption}}
\AtBeginDocument{%
\ifdefined\contentsname
  \renewcommand*\contentsname{Table of contents}
\else
  \newcommand\contentsname{Table of contents}
\fi
\ifdefined\listfigurename
  \renewcommand*\listfigurename{List of Figures}
\else
  \newcommand\listfigurename{List of Figures}
\fi
\ifdefined\listtablename
  \renewcommand*\listtablename{List of Tables}
\else
  \newcommand\listtablename{List of Tables}
\fi
\ifdefined\figurename
  \renewcommand*\figurename{Figure}
\else
  \newcommand\figurename{Figure}
\fi
\ifdefined\tablename
  \renewcommand*\tablename{Table}
\else
  \newcommand\tablename{Table}
\fi
}
\@ifpackageloaded{float}{}{\usepackage{float}}
\floatstyle{ruled}
\@ifundefined{c@chapter}{\newfloat{codelisting}{h}{lop}}{\newfloat{codelisting}{h}{lop}[chapter]}
\floatname{codelisting}{Listing}
\newcommand*\listoflistings{\listof{codelisting}{List of Listings}}
\makeatother
\makeatletter
\makeatother
\makeatletter
\@ifpackageloaded{caption}{}{\usepackage{caption}}
\@ifpackageloaded{subcaption}{}{\usepackage{subcaption}}
\makeatother
\usepackage{bookmark}
\IfFileExists{xurl.sty}{\usepackage{xurl}}{} % add URL line breaks if available
\urlstyle{same}
\hypersetup{
  pdftitle={Matrix Decomposition with Kronecker Weights},
  pdfauthor={Jan de Leeuw},
  colorlinks=true,
  linkcolor={blue},
  filecolor={Maroon},
  citecolor={Blue},
  urlcolor={Blue},
  pdfcreator={LaTeX via pandoc}}


\title{Matrix Decomposition with Kronecker Weights}
\author{Jan de Leeuw}
\date{April 24, 2026}
\begin{document}
\maketitle
\begin{abstract}
The familiar orthogonal decomposition of a matrix into row effect,
column effects, and interactions is generalized to a matrix space with
the Kronecker inner product.
\end{abstract}

\renewcommand*\contentsname{Table of contents}
{
\setcounter{tocdepth}{3}
\tableofcontents
}

\sectionbreak

\textbf{Note:} This is a working manuscript which will be
expanded/updated frequently. All suggestions for improvement are
welcome. All Rmd, tex, html, pdf, R, and C files are in the public
domain. Attribution will be appreciated, but is not required. The files
can be found at \url{https://github.com/deleeuw/decomp}

\sectionbreak

\section{Kronecker Weights}\label{kronecker-weights}

\(R\) and \(C\) are positive semidefinite matrices of order \(n\) and
\(m\). They define an inner product on the space
\(\mathbb{R}^{n\times m}\) of \(n\times m\) matrices as
\(\langle X, Y\rangle_{RC}=\text{tr}\ RXCY'\) and a corresponding
squared norm (pseudo-norm if either \(R\) or \(C\) is singular) as
\(\|X\|_{RC}^2=\langle X, X\rangle_{RC}\). From now on we leave out the
subscript \(RC\).

\section{Matrix Decomposition}\label{matrix-decomposition}

Define three subspaces of \(\mathbb{R}^{n\times m}\). We use \(e\) for
vectors with all elements equal to one. The number of elements in \(e\)
will be clear from the context, otherwise we will use \(e_m\) and
\(e_n\).

\begin{enumerate}
\def\labelenumi{\arabic{enumi}.}
\tightlist
\item
  \(\mathbb{L}_\mu\), constant matrices, \(X=\mu e_ne_m'\).
\item
  \(\mathbb{L}_a\), matrices with all columns the same, \(X=ae_m'\).
\item
  \(\mathbb{L}_b\), matrices with all rows the same \(X=e_nb'\).
\end{enumerate}

We will hoose \(a, b, \mu\) such that these three subspacs are
orthogonal in the Kronecker inner product. For this we need
\begin{subequations}
\begin{align}
\text{tr}\ R\mu ee'Cea'&=\mu a'Re\times e'Ce=0,\\
\text{tr}\ R\mu ee'Cbe'&=\mu b'Ce\times e'Re,\\
\text{tr}\ Rae'Cbe'&=a'Re\times b'Ce.
\end{align}
\end{subequations}

It follows the three subspaces are orthogonal if \(a'Re=b'Ce=0\). Thus
we redefine our subspaces as

\begin{enumerate}
\def\labelenumi{\arabic{enumi}.}
\tightlist
\item
  \(\mathbb{L}_\mu\): all \(X\) with \(X=\mu ee'\).
\item
  \(\mathbb{L}_a\): all \(X\) with \(X=ae'\) where \(a'Re=0\).
\item
  \(\mathbb{L}_b\): all \(X\) with \(X=eb'\) where \(b'Ce=0\).
\end{enumerate}

Now project, in the Kronecker metric, matrix \(X\) on each of these
orthogonal subspaces.

First on \(\mathbb{L}_\mu\). Minimize over \(\mu\) \begin{equation}
\|X-\mu ee'\|^2=\|X\|^2-2\mu e'RXCe+\mu^2e'Re\times e'Ce,
\end{equation} and thus \begin{equation}
\hat\mu=\frac{e'RXCe}{e'Re\times e'Ce}.\label{eq-mu}
\end{equation} Now over \(\mathbb{L}_a\). We have \begin{equation}
\|X-ae'\|^2=\|X\|^2-2a'RXCe+a'Ra\times e'Ce,
\end{equation} which must be minimized over all \(a'Re=0\). Use a
Lagrange multiplier \(\lambda\). Stationary equations are
\begin{equation}
RXCe-\lambda Re=e'Ce\times Ra.
\end{equation} Premultiply by \(e'\) gives \begin{equation}
\hat\lambda=\frac{e'RXCe}{e'Re},
\end{equation} and thus \begin{equation}
e'Ce\times\hat a=XCe-e\frac{e'RXCe}{e'Re},\label{eq-asolve}
\end{equation} or \begin{equation}
\hat a=\left(I-\frac{ee'R}{e'Re}\right)X\left(\frac{Ce}{e'Ce}\right).
\end{equation}

It is convenient to define the projectors \begin{align}
P:=\frac{ee'R}{e'Re},\label{eq-pdef}\\
Q:=\frac{Cee'}{e'Ce}.\label{eq-qdef}
\end{align} The projection on \(\mathbb{L}_a\) is \begin{equation}
\hat ae'=(I-P)XQ.
\end{equation}

Projecting on \(\mathbb{L}_b\) is similar, by symmetry. Thus
\begin{equation}
e\hat b'=PX(I-Q).
\end{equation} In addition, from \eqref{eq-mu}, \begin{equation}
\hat\mu ee'=PXQ.
\end{equation} It follows that if we define \begin{equation}
\hat H:=X-\hat\mu ee'-\hat ae'-e\hat b'\label{eq-hdef}
\end{equation} then \begin{equation}
\hat H=(I-P)X(I-Q).
\end{equation}

Define \(\mathbb{L}_h\) as the subspace of all matrices of this form, or
equivalently the matrices with \(HCe=0\) and \(e'RH=0\).

Because \((I-Q)Ce=0\) it follows that \begin{subequations}
\begin{equation}
\text{tr}\ R\hat HC(\hat\mu ee')=0
\end{equation}
and 
\begin{equation}
\text{tr}\ R\hat HC(e\hat b')=0
\end{equation}
and
\begin{equation}
\text{tr}\ R\hat ae'C\hat H'=0
\end{equation}
\end{subequations} Thus \(\mathbb{L}_h\) is also orthogonal in the
Kronecker inner product to the other three subspaces. It follows that in
the Kronecker norm \begin{equation}
\|X\|^2=\|PXQ\|^2+\|(I-P)XQ\|^2+\|PX(I-Q)\|^2+\|(I-P)X(I-Q)\|^2,
\end{equation} and thus we not only have a decomposition of the data
\(X\), but also of its Kronecker sum of squares.

It also follows that if \(Y=\mu ee'+ae'+eb'+H\), with \(a'Re=b'Ce=0\),
then \begin{multline}
\|X-Y\|^2=\|PXQ-\mu ee'\|^2+\|(I-P)XQ-ae'\|^2+\\\|PX(I-Q)-eb'\|^2+\|(I-P)X(I-Q)-H\|^2,
\end{multline} which can also be written as \begin{equation}
\|X-Y\|^2=\|(\hat\mu-\mu)ee'\|^2+\|(\hat a-a)e'\|^2+\|e(\hat b-b)'\|^2+\|\hat H-H\|^2.
\end{equation} Thus if we are fitting \(Y\) to \(X\) with separable
restrictions on \(a\), \(b\), and \(H\) (in addition to \(a'Re=b'Ce=0\))
then we can minimize each of the four components separately.

We have defined \(H\) in \eqref{eq-hdef} as the residual after
projecting on the three orthogonal spaces. The orthogonality of the
residual space is YAPT (Yet Another Pythogorean Theorem). Alternatively,
we could have defined \(\mathbb{L}_h\) as the space of \(n\times m\)
matrices satisfying \(HCe=0\) and \(H'Re=0\) and project on this space.
This will produce the same \(\hat H\). Here is a proof of this last
statement. The stationary equations for the projection, using two
vectors \(\phi\) and \(\psi\) of Lagrange multipliers are
\begin{subequations}
\begin{align}
H&=X-\phi e'-e\psi',\label{eq-hstat1}\\
HCe &= 0,\label{eq-hstat2}\\
H'Re&=0.\label{eq-hstat3}
\end{align}
\end{subequations} Postmultiply \eqref{eq-hstat1} by \(Ce\). Then we
must have \(0=XCe-\phi e'Ce-e\psi'Ce\) and thus \begin{subequations}
\begin{equation}
\phi=\frac{XCe}{e'Ce}-\frac{ee'C}{e'Ce}\psi.\label{eq-phicalc}
\end{equation}
In the same way
\begin{equation}
\psi=\frac{X'Re}{e'Re}-\frac{ee'R}{e'Re}\phi.\label{eq-psicalc}
\end{equation}
\end{subequations} Substitute \eqref{eq-psicalc} into
\eqref{eq-phicalc}. Then \begin{equation}
\phi=\frac{XCe}{e'Ce}-\frac{ee'C}{e'Ce}\left\{\frac{X'Re}{e'Re}-\frac{ee'R}{e'Re}\phi\right\}.
\end{equation} Collecting terms and using \eqref{eq-asolve} and
\eqref{eq-pdef} gives \((I-P)\phi=\hat a\). In the same way
\((I-Q)\psi=\hat b\). This implies that \(\phi=\hat a+\alpha e\) and
\(\psi=\hat b+\beta e\) for some scalars \(\alpha\) and \(\beta\), so
that \begin{equation}
\hat H=X-\hat ae'-e\hat b'-(\alpha+\beta)ee'.\label{eq-hcalc}
\end{equation} Premultiplying with \(e'R\) and postmultiplying with
\(Ce\), and then using \eqref{eq-hstat2} and \eqref{eq-hstat3}, gives
that \(\alpha+\beta=\hat\mu\). End of proof.

The developments in this paper can be generalized in several ways. In
the first place the same reasoning applies to multiway arrays. If the
array has \(p\) ``dimensions'' then there are \(2^p\) components.
Secondly, vectors \(e_n\) and \(e_m\) in the formulas can be replaced by
any vectors \(u\) and \(v\) with \(n\) and \(m\) elements. In fact, with
some additional notation, they can be replaced by matrices \(U\) and
\(V\). And finally the results apply to any finite sequence of
orthogonal subspaces in any (finite-dimensional) inner product space.

\sectionbreak

\section{Example}\label{example}

\begin{Shaded}
\begin{Highlighting}[]
\FunctionTok{set.seed}\NormalTok{(}\DecValTok{12345}\NormalTok{)}
\NormalTok{ymat }\OtherTok{\textless{}{-}} \FunctionTok{matrix}\NormalTok{(}\FunctionTok{rnorm}\NormalTok{(}\DecValTok{12}\NormalTok{), }\DecValTok{4}\NormalTok{, }\DecValTok{3}\NormalTok{)}
\NormalTok{rmat }\OtherTok{\textless{}{-}} \FunctionTok{crossprod}\NormalTok{(}\FunctionTok{matrix}\NormalTok{(}\FunctionTok{rnorm}\NormalTok{(}\DecValTok{40}\NormalTok{), }\DecValTok{10}\NormalTok{, }\DecValTok{4}\NormalTok{)) }\SpecialCharTok{/} \DecValTok{10}
\NormalTok{cmat }\OtherTok{\textless{}{-}} \FunctionTok{crossprod}\NormalTok{(}\FunctionTok{matrix}\NormalTok{(}\FunctionTok{rnorm}\NormalTok{(}\DecValTok{30}\NormalTok{), }\DecValTok{10}\NormalTok{, }\DecValTok{3}\NormalTok{)) }\SpecialCharTok{/} \DecValTok{10}
\end{Highlighting}
\end{Shaded}

\begin{verbatim}
[1] "ymat"
\end{verbatim}

\begin{verbatim}
           [,1]       [,2]       [,3]
[1,]  0.5855288  0.6058875 -0.2841597
[2,]  0.7094660 -1.8179560 -0.9193220
[3,] -0.1093033  0.6300986 -0.1162478
[4,] -0.4534972 -0.2761841  1.8173120
\end{verbatim}

\begin{verbatim}
[1] "cmat"
\end{verbatim}

\begin{verbatim}
           [,1]       [,2]       [,3]
[1,]  1.3744581  0.7610956 -0.2692483
[2,]  0.7610956  1.2428719 -0.1845269
[3,] -0.2692483 -0.1845269  1.3504256
\end{verbatim}

\begin{verbatim}
[1] "rmat"
\end{verbatim}

\begin{verbatim}
          [,1]       [,2]       [,3]       [,4]
[1,] 0.6606540  0.6316750  0.2060831  0.4388240
[2,] 0.6316750  1.5141559 -0.7485539  0.5940988
[3,] 0.2060831 -0.7485539  2.0102613 -0.4442019
[4,] 0.4388240  0.5940988 -0.4442019  1.3318363
\end{verbatim}

\sectionbreak

We now apply the decompose() function from the code section.

\begin{Shaded}
\begin{Highlighting}[]
\NormalTok{h}\OtherTok{\textless{}{-}}\FunctionTok{decompose}\NormalTok{(ymat, rmat, cmat)}
\end{Highlighting}
\end{Shaded}

The four components are:

\begin{verbatim}
[[1]]
            [,1]        [,2]        [,3]
[1,] -0.01383449 -0.01383449 -0.01383449
[2,] -0.01383449 -0.01383449 -0.01383449
[3,] -0.01383449 -0.01383449 -0.01383449
[4,] -0.01383449 -0.01383449 -0.01383449

[[2]]
          [,1]       [,2]      [,3]
[1,] 0.2414392 -0.3254692 0.1578904
[2,] 0.2414392 -0.3254692 0.1578904
[3,] 0.2414392 -0.3254692 0.1578904
[4,] 0.2414392 -0.3254692 0.1578904

[[3]]
            [,1]        [,2]        [,3]
[1,]  0.43727229  0.43727229  0.43727229
[2,] -0.59892198 -0.59892198 -0.59892198
[3,]  0.19675185  0.19675185  0.19675185
[4,]  0.07507512  0.07507512  0.07507512

[[4]]
            [,1]        [,2]       [,3]
[1,] -0.07934821  0.50791885 -0.8654880
[2,]  1.08078326 -0.87973030 -0.4644560
[3,] -0.53365990  0.77265039 -0.4570556
[4,] -0.75617703 -0.01195553  1.5981810
\end{verbatim}

\sectionbreak

The Kronecker inner products of the four components are:

\begin{verbatim}
              [,1]          [,2]          [,3]          [,4]
[1,]  6.027679e-03  5.637851e-18 -3.122502e-17 -6.938894e-18
[2,] -1.127570e-17  8.540491e-01 -1.734723e-16  1.110223e-16
[3,] -3.816392e-17  1.075529e-16  2.740579e+00 -1.110223e-16
[4,]  2.081668e-17 -5.551115e-17 -8.604228e-16  9.280312e+00
\end{verbatim}

The trace of the inner product matrix and the squared norm of the data:

\begin{verbatim}
[1] 12.88097 12.88097
\end{verbatim}

\sectionbreak

\section{Code}\label{code}

\begin{Shaded}
\begin{Highlighting}[]
\FunctionTok{set.seed}\NormalTok{(}\DecValTok{12345}\NormalTok{)}
\NormalTok{ymat }\OtherTok{\textless{}{-}} \FunctionTok{matrix}\NormalTok{(}\FunctionTok{rnorm}\NormalTok{(}\DecValTok{12}\NormalTok{), }\DecValTok{4}\NormalTok{, }\DecValTok{3}\NormalTok{)}
\NormalTok{rmat }\OtherTok{\textless{}{-}} \FunctionTok{crossprod}\NormalTok{(}\FunctionTok{matrix}\NormalTok{(}\FunctionTok{rnorm}\NormalTok{(}\DecValTok{40}\NormalTok{), }\DecValTok{10}\NormalTok{, }\DecValTok{4}\NormalTok{)) }\SpecialCharTok{/} \DecValTok{10}
\NormalTok{cmat }\OtherTok{\textless{}{-}} \FunctionTok{crossprod}\NormalTok{(}\FunctionTok{matrix}\NormalTok{(}\FunctionTok{rnorm}\NormalTok{(}\DecValTok{30}\NormalTok{), }\DecValTok{10}\NormalTok{, }\DecValTok{3}\NormalTok{)) }\SpecialCharTok{/} \DecValTok{10}

\NormalTok{decompose }\OtherTok{\textless{}{-}} \ControlFlowTok{function}\NormalTok{(xmat, rmat, cmat) \{}
\NormalTok{  n }\OtherTok{\textless{}{-}} \FunctionTok{nrow}\NormalTok{(xmat)}
\NormalTok{  m }\OtherTok{\textless{}{-}} \FunctionTok{ncol}\NormalTok{(xmat)}
\NormalTok{  pmat }\OtherTok{\textless{}{-}} \FunctionTok{matrix}\NormalTok{(}\FunctionTok{colSums}\NormalTok{(rmat), n, n, }\AttributeTok{byrow =} \ConstantTok{TRUE}\NormalTok{) }\SpecialCharTok{/} \FunctionTok{sum}\NormalTok{(rmat)}
\NormalTok{  qmat }\OtherTok{\textless{}{-}} \FunctionTok{matrix}\NormalTok{(}\FunctionTok{rowSums}\NormalTok{(cmat), m, m) }\SpecialCharTok{/} \FunctionTok{sum}\NormalTok{(cmat)}
\NormalTok{  px }\OtherTok{\textless{}{-}}\NormalTok{ pmat }\SpecialCharTok{\%*\%}\NormalTok{ xmat}
\NormalTok{  xq }\OtherTok{\textless{}{-}}\NormalTok{ xmat }\SpecialCharTok{\%*\%}\NormalTok{ qmat}
\NormalTok{  dd }\OtherTok{\textless{}{-}} \FunctionTok{rep}\NormalTok{(}\FunctionTok{list}\NormalTok{(}\DecValTok{0}\NormalTok{),}\DecValTok{4}\NormalTok{) }
\NormalTok{  dd[[}\DecValTok{1}\NormalTok{]] }\OtherTok{\textless{}{-}}\NormalTok{ px }\SpecialCharTok{\%*\%}\NormalTok{ qmat}
\NormalTok{  dd[[}\DecValTok{2}\NormalTok{]] }\OtherTok{\textless{}{-}}\NormalTok{ px }\SpecialCharTok{{-}}\NormalTok{ dd[[}\DecValTok{1}\NormalTok{]]}
\NormalTok{  dd[[}\DecValTok{3}\NormalTok{]] }\OtherTok{\textless{}{-}}\NormalTok{ xq }\SpecialCharTok{{-}}\NormalTok{ dd[[}\DecValTok{1}\NormalTok{]]}
\NormalTok{  dd[[}\DecValTok{4}\NormalTok{]] }\OtherTok{\textless{}{-}}\NormalTok{ xmat }\SpecialCharTok{{-}}\NormalTok{ dd[[}\DecValTok{1}\NormalTok{]] }\SpecialCharTok{{-}}\NormalTok{ dd[[}\DecValTok{2}\NormalTok{]] }\SpecialCharTok{{-}}\NormalTok{ dd[[}\DecValTok{3}\NormalTok{]]}
\NormalTok{  ip }\OtherTok{\textless{}{-}} \FunctionTok{matrix}\NormalTok{(}\DecValTok{0}\NormalTok{, }\DecValTok{4}\NormalTok{, }\DecValTok{4}\NormalTok{)}
  \ControlFlowTok{for}\NormalTok{ (i }\ControlFlowTok{in} \DecValTok{1}\SpecialCharTok{:}\DecValTok{4}\NormalTok{) \{}
    \ControlFlowTok{for}\NormalTok{ (j }\ControlFlowTok{in} \DecValTok{1}\SpecialCharTok{:}\DecValTok{4}\NormalTok{) \{}
\NormalTok{      ip[i, j] }\OtherTok{\textless{}{-}} \FunctionTok{kroneckerIP}\NormalTok{(dd[[i]], dd[[j]], rmat, cmat)}
\NormalTok{    \}}
\NormalTok{  \}}
\NormalTok{  ssq }\OtherTok{\textless{}{-}} \FunctionTok{c}\NormalTok{(}\FunctionTok{kroneckerIP}\NormalTok{(xmat, xmat, rmat, cmat), }\FunctionTok{sum}\NormalTok{(}\FunctionTok{diag}\NormalTok{(ip)))}
  \FunctionTok{return}\NormalTok{(}\FunctionTok{list}\NormalTok{(}\AttributeTok{dd =}\NormalTok{ dd, }\AttributeTok{ip =}\NormalTok{ ip, }\AttributeTok{ssq =}\NormalTok{ ssq))}
\NormalTok{\}}

\NormalTok{kroneckerIP }\OtherTok{\textless{}{-}} \ControlFlowTok{function}\NormalTok{(xmat, ymat, rmat, cmat) \{}
  \FunctionTok{return}\NormalTok{(}\FunctionTok{sum}\NormalTok{(ymat }\SpecialCharTok{*}\NormalTok{ (rmat }\SpecialCharTok{\%*\%}\NormalTok{ xmat }\SpecialCharTok{\%*\%}\NormalTok{ cmat)))}
\NormalTok{\}}
\end{Highlighting}
\end{Shaded}





\end{document}
